\subsection{User interface}
The main user-interface for snowman is the web-interface. It is based on
django's template-system for the generation of html which is displayed in the
web-browser on the user's computer.

The web-interface uses javascript to dynamicly be able to load/change the
content of the pages without refreshing the whole page. Important parts of the
webinterface is possible to download using AJAX, to update tables and simmilar.

Javascript is also used to improve the user-experience of the interface, due to
functionality which a moden computer-user expects from modern web-pages.

\subsubsection{Url's}

The first step of serving a page to a user is a HTTP request for a specific page
is recieved. The pages are uniqely identified by url's, like '/web/rules/'.

The url's are structured in a logical manner, where everything in the
webinterface is found under /web/, everything regarding the administration of
rules is found under /web/rules/, etc. The url's might contain dynamic data,
like pagenumbers (ie: /web/rules/rules/page/3/ which will give you the third
page of the rule-list), so that dynamic content can be loaded dynamicly when
needed.

Django uses the urls to identify which views the requested url corresponds to.
When a view is identified, the request (and its dynamic content) is passed to
the corresponding view.

\subsubsection{Views}

The views is the part of the user-interface which collects the required data
from the models-layer, and generates a response to the web client. The response
can contain a full web-page, or just a small subset of one if the request was
issued by AJAX.

The typical steps a view makes, to generate a response would be:
\begin{enumerate}
  \item Collect the requiered data from the database(s) using the \django{}
  models.
  \item Hand the collected data, and a template, to the template-engine to
  generate the HTML representing the webpage the user want.
  \item Send the created webpage as a response to the user.
\end{enumerate}

The view-system decides what to show, and then lets the template-system decide
how to show it. The view-system in itself is not creating any html or simmilar.
It soley creates data-structures.

If the request was for some data, and not a webpage, step 2 in the list above
might be skipped, and instead return the data as JSON or simmilar in step 3.

If errors oocures in the view-system, appropriate exception can be raised to
tell \django{} to return an error-code or an error-page to the user, instead of
what the user requested.

\subsubsection{Templates}

When the view-system have prepared all the data, the template-system takes over.
The job of the templates is to take the supplied data, and inject it into
pre-made skeletons of HTML. This way, all the design/display part of the
webinterface is collected to one single place, and that give us a very clear
line between data and design.
